\subsection{Cerebras Wafer-Scale Engine Architecture}

The WSE comprises a networked grid of independently executing compute cores (Processing Elements or PEs).
Each PE contains specialized message-handling infrastructure, which routes tagged 32-bit packets to neighboring PEs and/or to be processed locally, according to a programmed rule set.
Each PE is equipped with 48kB of private on-chip memory, which can be accessed within a single clock cycle;
communication to neighboring PEs, too, incurs low latency \cite{buitrago2021neocortex}.
PEs provide faculties for standard arithmetic and flow-control operations, as well as vectorized 32- and 16-bit integer and floating-point operations.
As Cerebras' driving market area is AI/ML, their product provides extensive integration frameworks for deep learning (e.g., TensorFlow).
However, they have also been active in applications to general-purpose HPC \cite{jacquelin2022scalable,rocki2020fast}, and have opened the full capabilities of the platform to outside developers through their Software Development Kit (SDK) \cite{selig2022cerebras}.
The WSE device is programmed by providing ``kernel'' code, written in the Cerebras Software Language (CSL), to be executed on each PE.
This language's programming model purposefully reflects underlying capabilities and particularities of the WSE architecture.
CSL organizes code within an event-driven framework, with the programmer defining tasks to be triggered in response to wavelet activation signals exchanged between --- and within --- PEs.
Scheduling of active tasks occurs via hardware-level microthreading, which allows some level of concurrency.
Special faculties are provided for asynchronous send-receive operations that exchange contiguous or temporally-allotted buffer data between PEs.
This Cerebras SDK includes a hardware emulator, providing means to compile and test CSL code from any Linux desktop environment, regardless of hardware access.
% Cerebras' current hardware offering is the CS-2 device, which offers 850,000 PEs and boasts a peak performance of 7.5 PetaFLOPs at 16-bit precision \cite{lie2023cerebras}.
% This model will shortly be succeeded by the CS-3, which implements nearly the same architectural model but with an increase of PE count to 900,000 and a few other spec bumps \cite{moore2024cerebras}.
% Notably, the CS-3 architecture boosts support for clustering, allowing interconnection of potentially thousands of WSE chips.
% Even with the current generation CS-2 hardware, HPC applications clustering 24 WSE chips to reach a PE count of 35 million have been reported \cite{ltaief2023scaling}.
